% ============================================================
%  Abstract
% ============================================================

\begin{abstract}
This paper presents Conceptual Topography (CT), a field-theoretic
account of conceptual landscapes, as the first application layer of
the VED$\times$IFGT framework.
A conceptual landscape is defined as the long-time geometry of the information
potential field $\Phi(x,\tau)$ under recursive token--field interaction:
a stable, sedimented structure that constrains the trajectories of thought
without being directly accessible to conscious inspection.

Within the VED$\times$IFGT framework, conceptual structure emerges from
three coupled processes.
First, the continuous information field $I(x,\tau)$, inherited from the
VED$\times$IFGT framework as a cognitive-window projection of
$I=f(1-C)$, generates an effective potential $\Phi = K * I$ that shapes
the passability landscape.
Second, discrete language tokens
$\ell_i = \mathcal{Q}[I,\Phi,C;\,\text{constraints}]$ act as
re-addressable discrete units that interact locally with $\nabla\Phi$,
deepening potential wells through repeated reuse.
Third, the recursive loop
$\Delta \to C \to I \to \Phi \to \JI \to \Delta$
drives irreversible sedimentation, producing the attractor geometry
$\Phi \seddto \Phistar$ that constitutes a mature conceptual landscape.

The theory yields precise structural definitions: a concept is a locally
stable region of $\Phi$ satisfying $\nabla\Phi(x) \approx 0$; understanding
is a temporary reduction in traversal cost across such a region; thought is
the trajectory of $\JI$ across the landscape.
The strict positivity $I(x,\tau) > 0$ excludes complete conceptual
closure as an ordinary dynamical state.
In nontrivial generated landscapes, residual asymmetry and boundary
structure allow further gradients and differentiations to reappear.

At the social scale, the same dynamics produce a shared potential geometry
maintained across agents through linguistic synchronization, corresponding to
the emergence of norms, institutions, and collective intelligence.

This work does not redefine any variable of the VED$\times$IFGT framework;
it applies the existing field-theoretic structure to the domain of cognition
and language.
\end{abstract}
